% Source-aligned English for printed p.215.
\clearpage
\SourcePage{215}

\LayerTarget{tohoku:comparison:equivariant-hom-stalk-to-fiber}
\LayerTarget{tohoku:equation:5.6.9}
\noindent bijective when ${\mathbf O}$ is a \textit{coherent} sheaf of left
\textit{Noetherian} rings and $A$ is coherent as an ${\mathbf O}$-module; and, under
these conditions on $X,G,{\mathbf O}$, the $U_x$-module $B(x)$ is injective whenever $B$
is an injective object of ${\mathbf C}^{{\mathbf O}(G)}$. (In the proof, the free ${\mathbf O}$-modules of
4.1 are to be replaced by $G$-${\mathbf O}$-modules of type $L(U)$ introduced in the
proof of Proposition 5.1.1.) It follows immediately that the \textit{natural
homomorphisms}

\LayerTarget{tohoku:comparison:equivariant-ext-stalk-to-fiber}
\LayerTarget{tohoku:equation:5.6.10}
\[
\tag{5.6.10}
\mathbf{Ext}_{{\mathbf O},G}^n(A,B)(y)
\to \operatorname{Ext}_{U_x}^n(A(x),B(x))
\]
\textit{deduced from homomorphisms} (5.6.9) \textit{are
isomorphisms if the group of homeomorphisms $G$ is assumed discontinuous in
the sense of \textup{5.3}, ${\mathbf O}$ is a coherent sheaf of left Noetherian rings, and $A$ is
a $G$-${\mathbf O}$-module coherent as an ${\mathbf O}$-module}. This statement contains both
Theorems 4.2.2 and 5.3.1 (the constant sheaf of rings $\mathbf Z$ being
coherent and Noetherian!) and, in the most important cases, specifies the
structure of the $\mathbf{Ext}_{{\mathbf O},G}^n(A,B)$.

\LayerTarget{tohoku:spectral-sequence:equivariant-ext-composite-functors}
From formulas (5.6.4) and (5.6.5), one can obtain as many spectral sequences
by applying Theorem 2.4.1. We shall be concerned only with those deduced from
(5.6.5), converging to $\operatorname{Ext}_{{\mathbf O},G}^{*}(X;A,B)$. It remains to
verify that, in each of these three formulas expressing $h_{{\mathbf O},G,A}$ as a
composite functor, the usual acyclicity condition is satisfied; this follows
from

\LayerTarget{tohoku:lemma:5.6.1}
\textsc{Lemma 5.6.1.}\quad \textit{Suppose that $B$ is an injective
object of ${\mathbf C}^{{\mathbf O}(G)}$. Then $\mathbf{Hom}_{{\mathbf O},G}(A,B)$ is a flasque sheaf on $Y$,
$\operatorname{Hom}_{\mathrm O}(A,B)$ is a $\Gamma^G$-acyclic $G$-module, and
$\mathbf{Hom}_{\mathrm O}(A,B)$ is a $\Gamma_X^G$-acyclic $G$-sheaf.}

\LayerTarget{tohoku:proof:lemma-5.6.1}
By Proposition 5.1.2, we are reduced to the case in which $B$ is the product
sheaf defined by a family $(B_x)_{x\in X}$ of injective $U_x$-modules $B_x$.
It follows that the abelian $G$-sheaf
$\mathbf H=\mathbf{Hom}_{\mathbf O}(A,B)$ is the product sheaf defined by the family of
$G_x$-modules $H_x=\operatorname{Hom}_{{\mathbf O}(x)}(A(x),B_x)$.
\LayerTarget{tohoku:construction:product-hom-sheaf}
Then $\mathbf{Hom}_{{\mathbf O},G}(A,B)$ is the product sheaf on $Y$ defined by the
family of groups
\LayerTarget{tohoku:equation:product-stalk-invariants}
$H_y=\left(\displaystyle\prod H_x\right)^G$ (the product ranging over the $x$ such that $f(x)=y$), and is therefore flasque. Assume for a moment that the $H_x$ are
$\Gamma^{G_x}$-acyclic $G_x$-modules (this will be proved in the corollary
below).
\LayerTarget{tohoku:construction:coinduced-hom-module}
Then the $G$-module $\overline H_x$ induced from the $G_x$-module $H_x$ by
contravariant extension of scalars is $\Gamma^G$-acyclic, as is well known;
hence
\LayerTarget{tohoku:equation:global-hom-product}
$\operatorname{Hom}_{\mathbf O}(A,B)=\Gamma_X\mathbf H
=\displaystyle\prod_{x\in X/G}\overline H_x
$ is $\Gamma^G$-acyclic. Finally, since $\mathbf H$ is flasque and hence $\Gamma_X$-acyclic, it follows
from the second spectral sequence of Theorem 5.2.1 that
\LayerTarget{tohoku:equation:equivariant-hom-cohomology-vanishing}
$H^n(X;G,\mathbf H)=H^n(G,\Gamma_X\mathbf H)$, which vanishes for $n>0$ by what we have seen. This completes the
proof of Lemma 5.6.1, provided that we prove the following special case:

\LayerTarget{tohoku:corollary:coinduced-hom-acyclicity}
\textsc{Corollary.}\quad \textit{Let $O$ be a ring with identity on
which a group $G$ acts, and let $U$ be the ring generated by $O$ and
$\mathbf Z(G)$, subject to the commutation relations
$g\lambda g^{-1}=\lambda^g$ for $g\in G$, $\lambda\in O$. Let $A$ and $B$ be
two $U$-modules, with $B$ injective. Then the $G$-module
$\operatorname{Hom}_{\mathbf O}(A,B)$ is $\Gamma^G$-acyclic. (One must not suppose that
it is injective!)}

\LayerTarget{tohoku:proof:coinduced-hom-acyclicity}
\LayerTarget{tohoku:definition:coinduced-module-b-prime}
Let $B'=\operatorname{Hom}_{Z}(\mathbf Z(G),B)$ be the group of maps
$f$ from $G$ to $B$; we

\TohokuContinuation{216}
